\Needspace{6\baselineskip}
\begin{center}\hypertarget{native-254-article182}{}\pdfbookmark[2]{182.}{article-182}182.\end{center}
Disquisitioni de gradu praecisionis incognitarum valoribus tribuendo,
quoties plures adsunt, praemittere oportebit considerationem accuratiorem
functionis $vv+v'v'+v''v''+\text{etc.}$, quam per $W$ denotabimus.

I. Statuamus $\tfrac12\frac{\mathrm{d}W}{\mathrm{d}p}=p'
=\lambda+\alpha p+\beta q+\gamma r+\delta s+\text{etc.}$, atque
$W-\frac{p'p'}{\alpha}=W'$, \SourceLeaf{255}patetque fieri $p'=P$, et,
quum sit $\frac{\mathrm{d}W'}{\mathrm{d}p}=\frac{\mathrm{d}W}{\mathrm{d}p}
-\frac{2p'}{\alpha}\cdot\frac{\mathrm{d}p'}{\mathrm{d}p}=0$, functionem
$W'$ a $p$ liberam fore. Coëfficiens $\alpha=aa+a'a'+a''a''+\text{etc.}$
manifesto semper erit quantitas positiva.

II. Perinde statuemus $\tfrac12\cdot\frac{\mathrm{d}W'}{\mathrm{d}q}=q'
=\lambda'+\beta'q+\gamma'r+\delta's+\text{etc.}$, atque
$W'-\frac{q'q'}{\beta'}=W''$, eritque
$q'=\tfrac12\cdot\frac{\mathrm{d}W}{\mathrm{d}q}
-\frac{p'}{\alpha}\cdot\frac{\mathrm{d}p'}{\mathrm{d}q}
=Q-\frac{\beta}{\alpha}\cdot p'$, atque
$\frac{\mathrm{d}W''}{\mathrm{d}q}=0$, unde patet, functionem $W''$ tum a
$p$ tum a $q$ liberam fore. Haec locum non haberent, si fieri posset
$\beta'=0$. Sed patet, $W'$ oriri ex $vv+v'v'+v''v''+\text{etc.}$,
eliminata quantitate $p$ ex $v$, $v'$, $v''$ etc. adiumento aequationis
$p'=0$; hinc $\beta'$ erit summa coëfficientium ipsius $qq$ in $vv$,
$v'v'$, $v''v''$ etc. post illam eliminationem; hi vero singuli
coëfficientes ipsi sunt quadrata, neque omnes simul evanescere possunt,
nisi in casu supra excluso, ubi incognitae indeterminatae manent. Patet
itaque, $\beta'$ esse debere quantitatem positivam.

III. Statuendo denuo $\tfrac12\cdot\frac{dW''}{dr}=r'
=\lambda''+\gamma''r+\delta''s+\text{etc.}$, atque
$W''-\frac{r'r'}{\gamma''}=W'''$, erit
$r'=R-\frac{\gamma}{\alpha}p'-\frac{\gamma'}{\beta'}q'$, atque $W'''$
libera tum a $p$, tum a $q$, tum a $r$. Ceterum coëfficientem $\gamma''$
necessario positivum fieri, simili modo probatur, ut in II. Facile
scilicet perspicitur, $\gamma''$ esse summam coëfficientium ipsius $rr$ in
$vv$, $v'v'$, $v''v''$ etc., postquam quantitates $p$ et $q$ adiumento
aequationum $p'=0$, $q'=0$ ex $v$, $v'$, $v''$ etc. eliminatae sunt.

IV. Eodem modo statuendo $\tfrac12\frac{dW'''}{ds}=s'
=\lambda'''+\delta'''s+\text{etc.}$,
$W^{\mathrm{IV}}=W'''-\frac{s's'}{\delta'''}$, erit
$s'=S-\frac{\delta}{\alpha}p'-\frac{\delta'}{\beta'}q'-\frac{\delta''}{\gamma''}r'$,
$W^{\mathrm{IV}}$ a $p$, $q$, $r$, $s$ libera, atque $\delta'''$ quantitas positiva.

V. Hoc modo, si praeter $p$, $q$, $r$, $s$ adhuc aliae incognitae adsunt,
ulterius progredi licebit, ita ut tandem habeatur
\[
W=\frac{1}{\alpha}p'p'+\frac{1}{\beta'}q'q'+\frac{1}{\gamma''}r'r'
+\frac{1}{\delta'''}s's'+\text{etc.}+\text{Const.},
\]
ubi omnes coëfficientes $\alpha$, $\beta'$, $\gamma''$, $\delta'''$ etc.
erunt quantitates positivae.

VI. Iam probabilitas alicuius systematis valorum determinatorum pro
quantitatibus $p$, $q$, $r$, $s$ etc. proportionalis\HistoricalNoteEight{Handschriftliche Bemerkung:] \emph{Aequalis} [producto] functioni[s\quad\raisebox{.65ex}{.}
\[
e^{-hh\left(\frac{1}{\alpha}p'p'+\frac{1}{\beta'}q'q'+\frac{1}{\gamma''}r'r'+\frac{1}{\delta'''}s's'+\text{etc.}\right)}
\]
per factorem constantem]
$\frac{h^\nu\BareRoot{(\alpha\beta'\gamma''\delta'''\ldots)}}{(\BareRoot{\pi})^\nu}e^{-hh\,\text{Const.}}$,
ubi Const. designat partem ultimam constantem ipsius $W$.}
est functioni $e^{-hhW}$; quamobrem,
\SourceLeaf{256}manente valore quantitatis $p$ indeterminato, probabilitas
systematis valorum determinatorum pro reliquis, proportionalis erit
integrali $\int e^{-hhW}\mathrm{d}p$, a $p=-\infty$ usque ad $p=+\infty$
extenso, quod per theorema ill. \textsc{Laplace} fit
\[
=h^{-1}\alpha^{-\frac12}\pi^{\frac12}
e^{-hh\left(\frac{1}{\beta'}q'q'+\frac{1}{\gamma''}r'r'+\frac{1}{\delta'''}s's'+\text{etc.}\right)};
\]
haecce itaque probabilitas proportionalis erit functioni $e^{-hhW'}$.
Perinde si insuper $q$ tamquam indeterminata tractatur, probabilitas
systematis valorum determinatorum pro $r$, $s$ etc. proportionalis erit
integrali $\int e^{-hhW'}\mathrm{d}q$, a $q=-\infty$ usque ad $q=+\infty$
extenso, quod fit
\[
=h^{-1}{\beta'}^{-\frac12}\pi^{\frac12}
e^{-hh\left(\frac{1}{\gamma''}r'r'+\frac{1}{\delta'''}s's'+\text{etc.}\right)},
\]
sive proportionalis functioni $e^{-hhW''}$. Prorsus simili modo, si etiam
$r$ tamquam indeterminata consideratur, probabilitas valorum
determinatorum pro reliquis $s$ etc. proportionalis erit functioni
$e^{-hhW'''}$ et sic porro. Supponamus, incognitarum numerum ad quatuor
ascendere; eadem enim conclusio valebit, si maior vel minor est. Valor
maxime probabilis ipsius $s$ hic erit
$=-\frac{\lambda'''}{\delta'''}\mathpunct{\raisebox{.45ex}{,}}$
probabilitasque, hunc a vero differentia $\sigma$ distare, proportionalis
erit functioni $e^{-hh\delta'''\sigma\sigma}$, unde concludimus, mensuram
praecisionis relativae isti determinationi tribuendae exprimi per
$\BareRoot{\delta'''}$, si mensura praecisionis observationibus primitivis
tribuendae statuatur $=1$.
